
Here is the full solution for a Timoshenko cantilever beam fixed at $x=0$ and loaded with a transverse force $F_z$ at the free end $x=L$.

The solution is based on two key kinematic variables: the transverse displacement $u_z(\reflenx)$ and the cross-section rotation $\theta_y(\reflenx)$. 
The sign convention used is:
\begin{itemize}
\item $F_z$ is a positive (upward) force.
\item $u_z$ is positive (upward) displacement.
\item $M_y$ is a positive moment (causing tension in the $-z$ fibers, i.e., ``happy face" bending).
\item $V_z$ is a positive (upward) shear force on a positive $\reflenx$-face.
\end{itemize}

\textit{Equilibrium}
First, we determine the internal shear force $V_z(\reflenx)$ and bending moment $M_y(\reflenx)$ by simple statics.

\begin{itemize}
\item
Shear Force $V_z(\reflenx)$:
    By summing forces on a free-body diagram cut at \(\reflenx\):
    $V_z(\reflenx) = F_z$
    The shear force is constant along the entire beam.
\item
Bending Moment $M_y(\reflenx)$:
    By summing moments about the cut at \(\reflenx\):
    \[M_y(\reflenx) + F_z(L-x) = 0
    \qquad\implies\qquad
    M_y(\reflenx) = -F_z(L-x)
    \]
    The moment is linear, varying from $M_y(0) = -F_z L$ at the root to $M_y(L) = 0$ at the free end.
\end{itemize}

\textit{Curvature and Shear Strain}
Next, we find the curvature and shear strain using the Timoshenko beam constitutive relations.



\textit{3. Rotation and Displacement}

Finally, we integrate the kinematic relationships and apply the boundary conditions at the fixed end (\(\reflenx=0\)):
\begin{itemize}
\item $u_z(0) = 0$ (no displacement)
\item $\theta_y(0) = 0$ (no rotation)
\end{itemize}
\begin{itemize}
\item 
Rotation $\theta_y(\reflenx)$:
    We integrate the curvature, $\kappa_y = \theta_y'$, from the fixed end:
    $\theta_y(\reflenx) = \int_{0}^{x} \kappa_y(\xi) d\xi + \theta_y(0)$
    $\theta_y(\reflenx) = \int_{0}^{x} -\frac{F_z(L-\xi)}{E I_y} d\xi = -\frac{F_z}{E I_y} \left[ L\xi - \frac{\xi^2}{2} \right]_{0}^{x}$
    $$\theta_y(\reflenx) = -\frac{F_z}{E I_y} \left( Lx - \frac{x^2}{2} \right)$$
\item
Displacement $u_z(\reflenx)$:
    We use the shear strain definition with
    $u_z'(\reflenx) = \frac{F_z}{G A_z} - \left[ -\frac{F_z}{E I_y} \left( Lx - \frac{x^2}{2} \right) \right] = \frac{F_z}{G A_z} + \frac{F_z}{E I_y} \left( Lx - \frac{x^2}{2} \right)$
    
    Now, we integrate $u_z'$ from the fixed end:
    $u_z(\reflenx) = \int_{0}^{x} u_z'(\xi) d\xi + u_z(0)$
    $u_z(\reflenx) = \int_{0}^{x} \left[ \frac{F_z}{G A_z} + \frac{F_z}{E I_y} \left( L\xi - \frac{\xi^2}{2} \right) \right] d\xi$
    $$
    u_z(\reflenx) = \left[ \frac{F_z \xi}{G A_z} + \frac{F_z}{E I_y} \left( \frac{L\xi^2}{2} - \frac{\xi^3}{6} \right) \right]_{0}^{x}
    $$
    
    This gives the final solution for displacement:
    $$
    u_z(\reflenx) 
    = \underbrace{\frac{F_z}{G A_z}\reflenx}_{\text{Shear deformation}}
    + \underbrace{\frac{F_z}{E I_y} \left( \frac{1}{2}L x^2 - \frac{1}{6} x^3 \right)}_{\text{Bending deformation}}
    $$
\end{itemize}

\paragraph{Summary of Full Solution}

\begin{itemize}
\item Strain measures are determined by inverting the constitutive relations. 
    The shear strain is directly proportional to the shear force:
    $\gamma_z(\reflenx) = \frac{V_z(\reflenx)}{G A_z}$
    $$\gamma_z(\reflenx) = \frac{F_z}{G A_z}$$
% \item
    The curvature is directly proportional to the bending moment:
    $\kappa_y(\reflenx) = \frac{M_y(\reflenx)}{E I_y}$
    $$\kappa_y(\reflenx) = -\frac{F_z(L-x)}{E I_y}$$
\item
Transverse Displacement $u_z(\reflenx)$:
    $$
    u_z(\reflenx) = \frac{F_z}{G A_z}\reflenx + \frac{F_z}{6 E I_y}\reflenx^2(3L-\reflenx)
    $$
\end{itemize}
